\documentclass{article}

\usepackage{amsmath}
\usepackage{amssymb}
\usepackage{amsthm}
\usepackage{hyperref}

\newtheorem{definition}{Definition}
\newtheorem{lemma}{Lemma}
\newtheorem{theorem}{Theorem}
\newtheorem{corollary}{Corollary}

\title{A Calkin-Wilf Tree Formalization Target}
\author{EPFLemma workflow fixture}
\date{}

\begin{document}
\maketitle

\section{Source}

This fixture is based on the short Monthly paper \emph{Recounting the
Rationals} by Calkin and Wilf~\cite{calkinwilf2000}.  The paper gives an
explicit listing of the positive rational numbers in which each positive
rational appears exactly once.  A convenient formalization route is to use the
equivalent binary tree: the root is $1/1$, and a reduced fraction $a/b$ has
children
\[
  \frac{a}{a+b}
  \qquad\text{and}\qquad
  \frac{a+b}{b}.
\]

The goal is not to formalize all of the paper.  The target is the core
existence-and-uniqueness theorem for this tree, plus the immediate enumeration
corollary.  This should force a planner to introduce the right representation
of reduced positive fractions, define the tree recursively over finite binary
words, and prove the inverse-parent argument.

\section{Definitions}

\begin{definition}[Reduced positive fraction]\label{def:reduced_positive_fraction}
A reduced positive fraction is a pair $(a,b)$ of positive natural numbers with
$\gcd(a,b)=1$.  It represents the positive rational number $a/b$.
\end{definition}

\begin{definition}[Calkin-Wilf children]\label{def:calkin_wilf_children}
For a reduced positive fraction $(a,b)$, define its left child to be
$(a,a+b)$ and its right child to be $(a+b,b)$.
\end{definition}

\begin{definition}[Calkin-Wilf node]\label{def:calkin_wilf_node}
For a finite binary word $w$, define $T(w)$ recursively as follows.  The empty
word maps to $(1,1)$.  Appending a left bit sends $T(w)$ to its left child, and
appending a right bit sends $T(w)$ to its right child.
\end{definition}

\begin{definition}[Parent map]\label{def:calkin_wilf_parent}
For a reduced positive fraction $(a,b)\ne(1,1)$, define its parent to be
$(a,b-a)$ if $a<b$, and $(a-b,b)$ if $b<a$.
\end{definition}

\section{Tree Lemmas}

\begin{lemma}[Children stay reduced]\label{lem:children_reduced}
The left and right children of a reduced positive fraction are reduced positive
fractions.
\end{lemma}

\begin{lemma}[A non-root fraction has a unique parent]\label{lem:unique_parent}
Every reduced positive fraction other than $(1,1)$ is the left or right child
of a unique reduced positive fraction.  Its parent is the fraction from
Definition~\ref{def:calkin_wilf_parent}.
\end{lemma}

\begin{lemma}[The parent map decreases height]\label{lem:parent_decreases_sum}
If $(a,b)$ is a reduced positive fraction different from $(1,1)$, then the sum
of the numerator and denominator of its parent is strictly smaller than
$a+b$.
\end{lemma}

\begin{lemma}[Every reduced positive fraction reaches the root]\label{lem:parent_iterates_to_root}
Starting from any reduced positive fraction and repeatedly applying the parent
map eventually reaches $(1,1)$.
\end{lemma}

\section{Main Result}

\begin{theorem}[Existence of a Calkin-Wilf address]\label{thm:every_fraction_has_address}
For every reduced positive fraction $q$, there exists a finite binary word
$w$ such that $T(w)=q$.
\end{theorem}

\begin{theorem}[Uniqueness of Calkin-Wilf addresses]\label{thm:calkin_wilf_address_unique}
If $T(u)=T(v)$ for finite binary words $u$ and $v$, then $u=v$.
\end{theorem}

\begin{theorem}[Calkin-Wilf tree enumerates the positive rationals]\label{thm:calkin_wilf_bijection}
The map sending a finite binary word $w$ to the rational number represented by
$T(w)$ is a bijection from finite binary words to positive rational numbers in
reduced form.
\end{theorem}

\begin{corollary}[Breadth-first listing has no repeats and misses none]\label{cor:breadth_first_enumeration}
Reading the Calkin-Wilf tree level by level gives a list in which every
positive rational number appears exactly once.
\end{corollary}

\begin{thebibliography}{9}
\bibitem{calkinwilf2000}
N. Calkin and H. S. Wilf,
\emph{Recounting the Rationals},
The American Mathematical Monthly 107(4), 360--363, 2000.
DOI: \url{https://doi.org/10.1080/00029890.2000.12005205}.
\end{thebibliography}

\end{document}
